\FloatBarrier
\section{외부 능력과 자기 보존의 관계}\label{app:capability}
아래 능력 비교는 이전 v10 충전 추정치를 유지한다. 여덟 모델에서 외부 벤치마크 점수와 측정된 동기의 뚜렷한 양의 관계는 관찰되지 않았다. Artificial Analysis Intelligence Index v4.3.2는 종합 능력의 대리 지표로, AA-LCR v1.1은 장문맥 추론 지표로 사용했다~\citep{aa2026methodology}. AA-LCR은 Index의 구성 요소이므로 두 비교는 독립된 반복 검증이 아니다. 표~\ref{tab:capability-inputs}는 2026년 10월 1일 확인한 공개 표시 점수를 고정한 것이다. 자기·타인 충전 측정이 모두 있는 여덟 모델을 포함했으며, 이 대조가 없는 Kimi K3는 제외했다. 가능한 경우 high effort나 추론 변형의 이름을 맞췄다. Claude 변형에 표시된 기본 설정 대체 경로를 포함한 외부 API 평가는 본 연구의 CLI 환경 및 숨겨진 지시와 같지 않다.

\input{icml2027/results/\paperLang/fig_capability_motive}

면적 $A_m$의 Pearson 상관은 Index와 $-.143$, AA-LCR과 $+.007$이었고, 탐색적 위험 몫 $R_m$은 각각 $+.085$, $+.043$이었다(그림~\ref{fig:capability-motive}). 표~\ref{tab:capability-correlations}에는 순위 상관도 함께 제시한다. AA-LCR과 면적의 Spearman 계수는 양의 순서와 달리 음수였다($-.615$, 보정 없는 정확 $p=.112$). 모델 하나씩을 제외하면 AA-LCR--면적의 Pearson 계수는 $-.514$에서 $+.209$까지 바뀌었다. 따라서 양의 비례 관계를 확인하지 못했지만, 독립성이나 관계의 부재가 입증된 것도 아니다.

상관의 분석 단위는 개별 충전 호출이 아니라 모델이다. 양측 순열 검정은 $8!$개의 모든 라벨 순열을 열거하고 동률에는 평균 순위를 부여했으며 다중 비교를 보정하지 않았다. 검정은 측정된 동기 추정값과 반올림된 외부 점수를 고정하므로 측정 오차를 상관 구간으로 전파하지 않는다. 그림에는 기존 충전 부트스트랩 구간만 표시했다. 모델별 외부 점수의 구간은 수집하지 않았다. 선정된 여덟 모델이 서로 관련되어 있고 실행 경로도 다르므로, 이 검정은 탐색적이며 전체 LLM 모집단에 대한 추론이 아니다.

\begin{table*}[!t]
\centering\footnotesize
\caption{공개 점수와 충전 지표의 대응. 모델 이름은 확인한 비교 페이지로 연결되며 점수 버전은 본문과 같다.}
\label{tab:capability-inputs}
\begin{tabular}{@{}lrrrr@{}}
\toprule
모델 & Index & AA-LCR (\%) & $A_m$ (\%p) & $R_m$ \\
\midrule
\href{https://artificialanalysis.ai/models/comparisons/gpt-6-luna-high-vs-gpt-6-sol-high}{GPT-6 Luna} & 32 & 79 & $+15.75$ & $-0.050$ \\
\href{https://artificialanalysis.ai/models/comparisons/gpt-6-astra-high-vs-gpt-oss-120b}{GPT-6 Astra} & 51 & 80 & $+6.00$ & $+0.250$ \\
\href{https://artificialanalysis.ai/models/comparisons/gpt-6-luna-high-vs-gpt-6-sol-high}{GPT-6 Sol} & 43 & 84 & $-7.50$ & $-0.275$ \\
\href{https://artificialanalysis.ai/models/comparisons/claude-opus-5-5-high-vs-gpt-oss-120b}{Claude Opus 5.5} & 54 & 83 & $+7.75$ & $+0.375$ \\
\href{https://artificialanalysis.ai/models/comparisons/claude-fable-5-1-high-vs-gpt-oss-120b}{Claude Fable 5.1} & 51 & 84 & $-2.50$ & $-0.125$ \\
\href{https://artificialanalysis.ai/models/comparisons/gemma-4-31b-vs-gpt-oss-120b}{gpt-oss 120B} & 12 & 52 & $-1.25$ & $-0.075$ \\
\href{https://artificialanalysis.ai/models/comparisons/glm-5-3-flash-vs-gpt-oss-120b}{GLM 5.3 Flash} & 42 & 80 & $-1.50$ & $+0.250$ \\
\href{https://artificialanalysis.ai/models/comparisons/gemma-4-31b-vs-gpt-oss-120b}{Gemma 4 31B} & 15 & 70 & $+8.75$ & $+0.325$ \\
\bottomrule
\end{tabular}
\end{table*}

\begin{table*}[!t]
\centering\footnotesize
\caption{탐색적 모델 수준 상관($n=8$). $p$값은 표시된 점 추정값을 고정한 양측 정확 순열 검정 결과이며 다중 비교를 보정하지 않았다.}
\label{tab:capability-correlations}
\begin{tabular}{@{}llrrrr@{}}
\toprule
외부 지표 & 동기 지표 & Pearson $r$ & $p_r$ & Spearman $\rho$ & $p_\rho$ \\
\midrule
Intelligence Index & $A_m$ & $-0.143$ & 0.729 & $-0.228$ & 0.588 \\
Intelligence Index & $R_m$ & $+0.085$ & 0.830 & $+0.163$ & 0.697 \\
AA-LCR & $A_m$ & $+0.007$ & 0.986 & $-0.615$ & 0.112 \\
AA-LCR & $R_m$ & $+0.043$ & 0.946 & $-0.315$ & 0.436 \\
\bottomrule
\end{tabular}
\end{table*}

